\documentclass[10pt,oneside]{article}
\usepackage[paperwidth=7in,paperheight=9.5in,top=.77in,bottom=.7in,left=.73in,right=.73in,headheight=15pt]{geometry}
\usepackage{fontspec,microtype,xcolor,hyperref,bookmark,listings,fancyhdr,titlesec,parskip,needspace}
\setmainfont{Liberation Serif}
\setsansfont{Liberation Sans}
\setmonofont[Scale=.88]{DejaVu Sans Mono}
\definecolor{ink}{HTML}{1D2D40}
\definecolor{accent}{HTML}{147B91}
\definecolor{muted}{HTML}{637587}
\definecolor{codebg}{HTML}{F2F5F8}
\definecolor{codekw}{HTML}{225BA0}
\definecolor{codestr}{HTML}{087A55}
\definecolor{codecom}{HTML}{718096}
\definecolor{outputbg}{HTML}{F8F9FA}
\definecolor{rule}{HTML}{D8E1E9}
\hypersetup{colorlinks=true,linkcolor=accent,urlcolor=accent,pdftitle={Go by Example - Edited Go Reference},pdfauthor={Original: Mark McGranaghan and contributors; edited reference adaptation},bookmarksnumbered=true,pdfstartview=FitH}
\titleformat{\section}{\sffamily\Large\bfseries\color{ink}}{\thesection}{.75em}{}
\titlespacing*{\section}{0pt}{2pt}{13pt}
\titleformat{\subsection}{\sffamily\small\bfseries\color{accent}}{}{0pt}{}
\titlespacing*{\subsection}{0pt}{10pt}{5pt}
\setlength{\parindent}{0pt}
\setlength{\parskip}{6pt plus 1pt}
\setcounter{tocdepth}{1}
\renewcommand{\contentsname}{Contents}
\pagestyle{fancy}
\fancyhf{}
\fancyhead[L]{\sffamily\scriptsize\bfseries\color{accent} GO BY EXAMPLE}
\fancyhead[R]{\sffamily\scriptsize\color{muted} EDITED GO REFERENCE}
\fancyfoot[L]{}
\fancyfoot[R]{\sffamily\small\color{ink}\thepage}
\renewcommand{\headrulewidth}{.35pt}
\renewcommand{\footrulewidth}{0pt}
\lstdefinelanguage{Go}{
  morekeywords={package,import,func,return,var,const,type,struct,interface,map,chan,go,defer,if,else,for,range,switch,case,default,break,continue,select,fallthrough},
  sensitive=true,morecomment=[l]{//},morecomment=[s]{/*}{*/},morestring=[b]",morestring=[b]`,
}
\lstdefinestyle{gocode}{language=Go,basicstyle=\color{ink}\ttfamily\small,keywordstyle=\color{codekw}\bfseries,
  stringstyle=\color{codestr},commentstyle=\color{codecom}\itshape,
  backgroundcolor=\color{codebg},frame=single,framerule=0pt,
  xleftmargin=0pt,xrightmargin=0pt,framexleftmargin=9pt,framexrightmargin=9pt,
  framextopmargin=7pt,framexbottommargin=7pt,
  aboveskip=5pt,belowskip=5pt,columns=fullflexible,keepspaces=true,showstringspaces=false,
  breaklines=true,breakatwhitespace=true,tabsize=4,upquote=true,
  literate={->}{{->}}2}
\lstdefinestyle{terminal}{basicstyle=\color{ink}\ttfamily\small,backgroundcolor=\color{outputbg},
  xleftmargin=0pt,xrightmargin=0pt,frame=single,framerule=0pt,
  framexleftmargin=9pt,framexrightmargin=9pt,framextopmargin=5pt,framexbottommargin=5pt,
  breaklines=true,columns=fullflexible,keepspaces=true,showstringspaces=false,aboveskip=3pt,belowskip=6pt}
\newcommand{\sectionline}{\par\vspace{4pt}{\color{rule}\hrule height .5pt}\vspace{7pt}\color{ink}\noindent}
\begin{document}
\color{ink}
\begin{titlepage}
\thispagestyle{empty}
\vspace*{1.0in}
{\sffamily\Huge\bfseries\color{ink}Go by Example\par}
\vspace{16pt}
{\sffamily\LARGE\color{accent}The Edited Reference\par}
\vspace{25pt}
\sectionline
\vspace{12pt}
{\Large\color{ink}Complete and self-contained Go examples.\par}
\vspace{16pt}
{\large\raggedright The original Go by Example project is by Mark McGranaghan and contributors, including Eli Bendersky. The source material is available at gobyexample.com and github.com/mmcgrana/gobyexample.\par}
\vfill
{\sffamily\small\color{muted}Adapted, corrected, and reformatted from the Go by Example project.\\Attribution and license appear in the opening pages.\par}
\end{titlepage}
\pagenumbering{roman}
\tableofcontents
\clearpage
\section*{How to use this reference}
Each chapter begins on a new page. The single \textbf{Example} listing in each chapter is complete Go source rather than a disconnected fragment. Save it as \texttt{main.go} and run \texttt{go run main.go}, except where a chapter explicitly describes a different file layout or a persistent server.

\textbf{Every chapter has an output example} directly below its code. For predictable programs, the listing shows the expected output. When results depend on goroutine scheduling, elapsed time, command-line arguments, the operating system, files, a remote service, or an external request, the listing shows \emph{one representative run}; a brief note explains what may vary. Terminal output can combine standard output and standard error when stated. A modern Go compiler is recommended; the iterator chapter requires Go 1.23 or later.

\textbf{Quick reading guide.} Declarations at package scope appear before \texttt{main}. Imports appear only when needed. The point of a chapter is the rule immediately above its example; the program supplies the runnable context needed to understand that rule.
\clearpage
\section*{Attribution and license}
The original Go by Example project is by Mark McGranaghan and contributors, including Eli Bendersky. The source material is available at \href{https://gobyexample.com}{gobyexample.com} and \href{https://github.com/mmcgrana/gobyexample}{github.com/mmcgrana/gobyexample}.

The original work is licensed under \href{https://creativecommons.org/licenses/by/3.0/}{Creative Commons Attribution 3.0 Unported (CC BY 3.0)}. This edition is an adaptation: the presentation, explanatory text, and sample programs have been revised substantially to make the material useful as a compact reference. The original authors do not endorse this adaptation.
\clearpage
\pagenumbering{arabic}
\clearpage
\section{Hello World}
\sectionline
A Go program starts in package main and executes func main. The fmt package prints formatted values.\par
\subsection*{Example}
\begin{lstlisting}[style=gocode]
package main

import "fmt"

func main() {
	fmt.Println("hello world")
}
\end{lstlisting}
\subsection*{Output}
\begin{lstlisting}[style=terminal]
hello world
\end{lstlisting}
\clearpage
\section{Values}
\sectionline
Go has strings, integers, floating-point values, and booleans. Operators act on compatible types.\par
\subsection*{Example}
\begin{lstlisting}[style=gocode]
package main

import "fmt"

func main() {
	fmt.Println("go" + "lang")
	fmt.Println(1 + 1)
	fmt.Println(7.0 / 3.0)
	fmt.Println(true && false)
	fmt.Println(true || false)
	fmt.Println(!true)
}
\end{lstlisting}
\subsection*{Output}
\begin{lstlisting}[style=terminal]
golang
2
2.3333333333333335
false
true
false
\end{lstlisting}
\clearpage
\section{Variables}
\sectionline
Use var for an explicit declaration and := for a short declaration inside a function. Uninitialized variables receive their zero value.\par
\subsection*{Example}
\begin{lstlisting}[style=gocode]
package main

import "fmt"

func main() {
	var name string = "Go"
	count := 3
	var ready bool
	fmt.Println(name, count, ready)
	count = 4
	fmt.Println(count)
}
\end{lstlisting}
\subsection*{Output}
\begin{lstlisting}[style=terminal]
Go 3 false
4
\end{lstlisting}
\clearpage
\section{Constants}
\sectionline
Constants are immutable. Untyped numeric constants can retain precision until a context gives them a type.\par
\subsection*{Example}
\begin{lstlisting}[style=gocode]
package main

import (
	"fmt"
	"math"
)

const message = "constant"
const big = 3e10

func main() {
	fmt.Println(message)
	fmt.Println(int64(big))
	fmt.Printf("%.2f\n", math.Sin(math.Pi/2))
}
\end{lstlisting}
\subsection*{Output}
\begin{lstlisting}[style=terminal]
constant
30000000000
1.00
\end{lstlisting}
\clearpage
\section{For}
\sectionline
for covers counted iteration, condition-only loops, and iteration over collections. break exits a loop; continue skips to the next iteration.\par
\subsection*{Example}
\begin{lstlisting}[style=gocode]
package main

import "fmt"

func main() {
	for i := 0; i < 3; i++ {
		fmt.Println(i)
	}
	n := 3
	for n > 0 {
		n--
	}
	fmt.Println("done")
}
\end{lstlisting}
\subsection*{Output}
\begin{lstlisting}[style=terminal]
0
1
2
done
\end{lstlisting}
\clearpage
\section{If/Else}
\sectionline
if conditions need no parentheses, but braces are mandatory. A short statement before the condition scopes a variable to the entire if/else chain.\par
\subsection*{Example}
\begin{lstlisting}[style=gocode]
package main

import "fmt"

func main() {
	if n := 9; n < 0 {
		fmt.Println("negative")
	} else if n < 10 {
		fmt.Println("one digit")
	} else {
		fmt.Println("many digits")
	}
}
\end{lstlisting}
\subsection*{Output}
\begin{lstlisting}[style=terminal]
one digit
\end{lstlisting}
\clearpage
\section{Switch}
\sectionline
switch chooses a branch without implicit fallthrough. Multiple values can share a case; a switch with no expression acts like a chain of conditions.\par
\subsection*{Example}
\begin{lstlisting}[style=gocode]
package main

import "fmt"

func main() {
	day := "Saturday"
	switch day {
	case "Saturday", "Sunday":
		fmt.Println("weekend")
	default:
		fmt.Println("weekday")
	}
}
\end{lstlisting}
\subsection*{Output}
\begin{lstlisting}[style=terminal]
weekend
\end{lstlisting}
\clearpage
\section{Arrays}
\sectionline
Arrays have fixed length, and that length is part of their type. Indexes start at zero; assigning an array copies its elements.\par
\subsection*{Example}
\begin{lstlisting}[style=gocode]
package main

import "fmt"

func main() {
	a := [3]int{1, 2, 3}
	b := a
	b[0] = 99
	fmt.Println(a)
	fmt.Println(b)
	fmt.Println(len(a))
}
\end{lstlisting}
\subsection*{Output}
\begin{lstlisting}[style=terminal]
[1 2 3]
[99 2 3]
3
\end{lstlisting}
\clearpage
\section{Slices}
\sectionline
Slices are dynamically sized views over arrays. append may allocate new storage, so use its return value.\par
\subsection*{Example}
\begin{lstlisting}[style=gocode]
package main

import "fmt"

func main() {
	a := []int{1, 2, 3}
	b := a[1:]
	b[0] = 20
	a = append(a, 4)
	fmt.Println(a)
	fmt.Println(len(a), cap(a))
}
\end{lstlisting}
\subsection*{Output}
\begin{lstlisting}[style=terminal]
[1 20 3 4]
4 6
\end{lstlisting}
\clearpage
\section{Maps}
\sectionline
Maps associate comparable keys with values. A lookup can return both the value and a boolean indicating whether the key exists.\par
\subsection*{Example}
\begin{lstlisting}[style=gocode]
package main

import "fmt"

func main() {
	ages := map[string]int{"Ada": 36}
	ages["Bob"] = 41
	age, ok := ages["Ada"]
	fmt.Println(age, ok)
	delete(ages, "Bob")
	_, ok = ages["Bob"]
	fmt.Println(ok)
}
\end{lstlisting}
\subsection*{Output}
\begin{lstlisting}[style=terminal]
36 true
false
\end{lstlisting}
\clearpage
\section{Functions}
\sectionline
Functions name their parameters and result types. A function can return a value or return nothing.\par
\subsection*{Example}
\begin{lstlisting}[style=gocode]
package main

import "fmt"

func add(a, b int) int {
	return a + b
}

func main() {
	fmt.Println(add(2, 3))
}
\end{lstlisting}
\subsection*{Output}
\begin{lstlisting}[style=terminal]
5
\end{lstlisting}
\clearpage
\section{Multiple Return Values}
\sectionline
A function may return multiple results. A common pattern returns a value and an error.\par
\subsection*{Example}
\begin{lstlisting}[style=gocode]
package main

import "fmt"

func divide(a, b int) (int, bool) {
	if b == 0 {
		return 0, false
	}
	return a / b, true
}

func main() {
	n, ok := divide(12, 3)
	fmt.Println(n, ok)
}
\end{lstlisting}
\subsection*{Output}
\begin{lstlisting}[style=terminal]
4 true
\end{lstlisting}
\clearpage
\section{Variadic Functions}
\sectionline
A variadic parameter accepts zero or more arguments. Inside the function, it behaves as a slice.\par
\subsection*{Example}
\begin{lstlisting}[style=gocode]
package main

import "fmt"

func sum(nums ...int) int {
	total := 0
	for _, n := range nums {
		total += n
	}
	return total
}

func main() {
	values := []int{1, 2, 3}
	fmt.Println(sum(values...))
}
\end{lstlisting}
\subsection*{Output}
\begin{lstlisting}[style=terminal]
6
\end{lstlisting}
\clearpage
\section{Closures}
\sectionline
A closure retains access to variables in its surrounding function, even after that outer function returns.\par
\subsection*{Example}
\begin{lstlisting}[style=gocode]
package main

import "fmt"

func counter() func() int {
	n := 0
	return func() int {
		n++
		return n
	}
}

func main() {
	next := counter()
	fmt.Println(next(), next(), next())
}
\end{lstlisting}
\subsection*{Output}
\begin{lstlisting}[style=terminal]
1 2 3
\end{lstlisting}
\clearpage
\section{Recursion}
\sectionline
Recursive functions call themselves and must have a base case to terminate.\par
\subsection*{Example}
\begin{lstlisting}[style=gocode]
package main

import "fmt"

func factorial(n int) int {
	if n <= 1 {
		return 1
	}
	return n * factorial(n-1)
}

func main() {
	fmt.Println(factorial(5))
}
\end{lstlisting}
\subsection*{Output}
\begin{lstlisting}[style=terminal]
120
\end{lstlisting}
\clearpage
\section{Range over Built-in Types}
\sectionline
range iterates over slices, arrays, strings, maps, and other iterable values. An underscore discards an unused index or value.\par
\subsection*{Example}
\begin{lstlisting}[style=gocode]
package main

import "fmt"

func main() {
	nums := []int{2, 3, 4}
	total := 0
	for _, n := range nums {
		total += n
	}
	fmt.Println(total)
	for i, n := range nums {
		fmt.Printf("%d:%d ", i, n)
	}
	fmt.Println()
}
\end{lstlisting}
\subsection*{Output}
\begin{lstlisting}[style=terminal]
9
0:2 1:3 2:4 
\end{lstlisting}
\clearpage
\section{Pointers}
\sectionline
A pointer holds the address of a variable. Passing a pointer allows a function to modify the original value.\par
\subsection*{Example}
\begin{lstlisting}[style=gocode]
package main

import "fmt"

func zero(p *int) {
	*p = 0
}

func main() {
	n := 7
	zero(&n)
	fmt.Println(n)
}
\end{lstlisting}
\subsection*{Output}
\begin{lstlisting}[style=terminal]
0
\end{lstlisting}
\clearpage
\section{Strings and Runes}
\sectionline
A string stores bytes, usually UTF-8 encoded text. range over a string decodes Unicode code points (runes).\par
\subsection*{Example}
\begin{lstlisting}[style=gocode]
package main

import (
	"fmt"
	"unicode/utf8"
)

func main() {
	text := "Go! 世界"
	fmt.Println(len(text))
	fmt.Println(utf8.RuneCountInString(text))
	for _, r := range "Go" {
		fmt.Printf("%c ", r)
	}
	fmt.Println()
}
\end{lstlisting}
\subsection*{Output}
\begin{lstlisting}[style=terminal]
10
6
G o 
\end{lstlisting}
\clearpage
\section{Structs}
\sectionline
Structs group fields into a named type. Construct them using field names to avoid order-dependent initialization.\par
\subsection*{Example}
\begin{lstlisting}[style=gocode]
package main

import "fmt"

type Person struct {
	Name string
	Age  int
}

func main() {
	p := Person{Name: "Ada", Age: 36}
	p.Age++
	fmt.Println(p.Name, p.Age)
}
\end{lstlisting}
\subsection*{Output}
\begin{lstlisting}[style=terminal]
Ada 37
\end{lstlisting}
\clearpage
\section{Methods}
\sectionline
A method is a function with a receiver. Pointer receivers can modify a value and avoid copying large structs.\par
\subsection*{Example}
\begin{lstlisting}[style=gocode]
package main

import "fmt"

type Rectangle struct {
	Width, Height int
}

func (r Rectangle) Area() int {
	return r.Width * r.Height
}

func (r *Rectangle) Scale(n int) {
	r.Width *= n
	r.Height *= n
}

func main() {
	r := Rectangle{3, 4}
	r.Scale(2)
	fmt.Println(r.Area())
}
\end{lstlisting}
\subsection*{Output}
\begin{lstlisting}[style=terminal]
48
\end{lstlisting}
\clearpage
\section{Interfaces}
\sectionline
An interface describes behavior via a method set. A concrete type satisfies an interface implicitly by providing those methods.\par
\subsection*{Example}
\begin{lstlisting}[style=gocode]
package main

import "fmt"

type Speaker interface {
	Speak() string
}

type Dog struct{}

func (Dog) Speak() string { return "woof" }

func announce(s Speaker) {
	fmt.Println(s.Speak())
}

func main() {
	announce(Dog{})
}
\end{lstlisting}
\subsection*{Output}
\begin{lstlisting}[style=terminal]
woof
\end{lstlisting}
\clearpage
\section{Enums}
\sectionline
Go has no built-in enum declaration. Named types plus iota provide readable sets of related constants.\par
\subsection*{Example}
\begin{lstlisting}[style=gocode]
package main

import "fmt"

type State int

const (
	Idle State = iota
	Running
	Stopped
)

func main() {
	s := Running
	fmt.Println(s == Running, s)
}
\end{lstlisting}
\subsection*{Output}
\begin{lstlisting}[style=terminal]
true 1
\end{lstlisting}
\clearpage
\section{Struct Embedding}
\sectionline
Embedding promotes fields and methods of an embedded type while retaining normal composition.\par
\subsection*{Example}
\begin{lstlisting}[style=gocode]
package main

import "fmt"

type Base struct{ ID int }

func (b Base) Label() string { return fmt.Sprint(b.ID) }

type Job struct {
	Base
	Name string
}

func main() {
	j := Job{Base: Base{ID: 7}, Name: "build"}
	fmt.Println(j.ID, j.Label(), j.Name)
}
\end{lstlisting}
\subsection*{Output}
\begin{lstlisting}[style=terminal]
7 7 build
\end{lstlisting}
\clearpage
\section{Generics}
\sectionline
Type parameters make functions work with multiple types while preserving static type checking. Constraints restrict permitted operations.\par
\subsection*{Example}
\begin{lstlisting}[style=gocode]
package main

import "fmt"

func Index[T comparable](items []T, want T) int {
	for i, v := range items {
		if v == want {
			return i
		}
	}
	return -1
}

func main() {
	fmt.Println(Index([]string{"a", "b"}, "b"))
}
\end{lstlisting}
\subsection*{Output}
\begin{lstlisting}[style=terminal]
1
\end{lstlisting}
\clearpage
\section{Range over Iterators}
\sectionline
Iterator functions can yield a sequence to range. The yield callback returns false when the consumer stops early (Go 1.23+).\par
\subsection*{Example}
\begin{lstlisting}[style=gocode]
package main

import "fmt"

func Count(n int) func(func(int) bool) {
	return func(yield func(int) bool) {
		for i := 0; i < n; i++ {
			if !yield(i) {
				return
			}
		}
	}
}

func main() {
	for n := range Count(3) {
		fmt.Println(n)
	}
}
\end{lstlisting}
\subsection*{Output}
\begin{lstlisting}[style=terminal]
0
1
2
\end{lstlisting}
\clearpage
\section{Errors}
\sectionline
Go errors are ordinary values. Check err explicitly and add useful context with fmt.Errorf and \%w.\par
\subsection*{Example}
\begin{lstlisting}[style=gocode]
package main

import (
	"errors"
	"fmt"
)

func load(name string) error {
	if name == "" {
		return errors.New("empty name")
	}
	return nil
}

func main() {
	if err := load(""); err != nil {
		fmt.Println("load failed:", err)
	}
}
\end{lstlisting}
\subsection*{Output}
\begin{lstlisting}[style=terminal]
load failed: empty name
\end{lstlisting}
\clearpage
\section{Custom Errors}
\sectionline
Custom error types can carry structured context. errors.As finds a matching error type through wrapping.\par
\subsection*{Example}
\begin{lstlisting}[style=gocode]
package main

import (
	"errors"
	"fmt"
)

type StatusError struct{ Code int }

func (e *StatusError) Error() string {
	return fmt.Sprintf("status %d", e.Code)
}

func main() {
	err := fmt.Errorf("request: %w", &StatusError{404})
	var target *StatusError
	if errors.As(err, &target) {
		fmt.Println(target.Code)
	}
}
\end{lstlisting}
\subsection*{Output}
\begin{lstlisting}[style=terminal]
404
\end{lstlisting}
\clearpage
\section{Goroutines}
\sectionline
A goroutine runs a function concurrently. Synchronize completion explicitly; sleeping is not a reliable synchronization technique.\par
\subsection*{Example}
\begin{lstlisting}[style=gocode]
package main

import (
	"fmt"
	"sync"
)

func main() {
	var wg sync.WaitGroup
	wg.Add(1)
	go func() {
		defer wg.Done()
		fmt.Println("from goroutine")
	}()
	wg.Wait()
}
\end{lstlisting}
\subsection*{Output}
\begin{lstlisting}[style=terminal]
from goroutine
\end{lstlisting}
\clearpage
\section{Channels}
\sectionline
Channels transfer typed values between goroutines. An unbuffered channel pairs a sender with a receiver.\par
\subsection*{Example}
\begin{lstlisting}[style=gocode]
package main

import "fmt"

func main() {
	ch := make(chan string)
	go func() { ch <- "hello" }()
	fmt.Println(<-ch)
}
\end{lstlisting}
\subsection*{Output}
\begin{lstlisting}[style=terminal]
hello
\end{lstlisting}
\clearpage
\section{Channel Buffering}
\sectionline
Buffered channels hold a limited number of queued values. A send blocks only when the buffer is full.\par
\subsection*{Example}
\begin{lstlisting}[style=gocode]
package main

import "fmt"

func main() {
	ch := make(chan int, 2)
	ch <- 10
	ch <- 20
	fmt.Println(<-ch, <-ch)
}
\end{lstlisting}
\subsection*{Output}
\begin{lstlisting}[style=terminal]
10 20
\end{lstlisting}
\clearpage
\section{Channel Synchronization}
\sectionline
A channel can signal that work finished. Closing a channel is a common broadcast signal when no more values will arrive.\par
\subsection*{Example}
\begin{lstlisting}[style=gocode]
package main

import "fmt"

func main() {
	done := make(chan struct{})
	go func() {
		fmt.Println("working")
		close(done)
	}()
	<-done
	fmt.Println("finished")
}
\end{lstlisting}
\subsection*{Output}
\begin{lstlisting}[style=terminal]
working
finished
\end{lstlisting}
\clearpage
\section{Channel Directions}
\sectionline
Channel direction annotations express intent and let the compiler reject mistaken sends or receives.\par
\subsection*{Example}
\begin{lstlisting}[style=gocode]
package main

import "fmt"

func send(ch chan<- string) {
	ch <- "ping"
}

func receive(ch <-chan string) {
	fmt.Println(<-ch)
}

func main() {
	ch := make(chan string)
	go send(ch)
	receive(ch)
}
\end{lstlisting}
\subsection*{Output}
\begin{lstlisting}[style=terminal]
ping
\end{lstlisting}
\clearpage
\section{Select}
\sectionline
select waits for one ready channel operation. A default branch makes the select non-blocking.\par
\subsection*{Example}
\begin{lstlisting}[style=gocode]
package main

import "fmt"

func main() {
	ch := make(chan int, 1)
	ch <- 42
	select {
	case n := <-ch:
		fmt.Println(n)
	default:
		fmt.Println("not ready")
	}
}
\end{lstlisting}
\subsection*{Output}
\begin{lstlisting}[style=terminal]
42
\end{lstlisting}
\clearpage
\section{Timeouts}
\sectionline
Use context deadlines to bound operations instead of waiting forever. Always call cancel to release associated resources.\par
\subsection*{Example}
\begin{lstlisting}[style=gocode]
package main

import (
	"context"
	"fmt"
	"time"
)

func main() {
	ctx, cancel := context.WithTimeout(context.Background(), time.Millisecond)
	defer cancel()
	<-ctx.Done()
	fmt.Println(ctx.Err())
}
\end{lstlisting}
\subsection*{Output}
\begin{lstlisting}[style=terminal]
context deadline exceeded
\end{lstlisting}
\clearpage
\section{Non-Blocking Channel Operations}
\sectionline
A select with default probes channels without waiting. That is useful for optional work, but not as a substitute for proper synchronization.\par
\subsection*{Example}
\begin{lstlisting}[style=gocode]
package main

import "fmt"

func main() {
	ch := make(chan string)
	select {
	case msg := <-ch:
		fmt.Println(msg)
	default:
		fmt.Println("no message")
	}
}
\end{lstlisting}
\subsection*{Output}
\begin{lstlisting}[style=terminal]
no message
\end{lstlisting}
\clearpage
\section{Closing Channels}
\sectionline
Only the sender should normally close a channel. A receiver can test whether a receive succeeded with the two-result form.\par
\subsection*{Example}
\begin{lstlisting}[style=gocode]
package main

import "fmt"

func main() {
	jobs := make(chan int, 2)
	jobs <- 1
	jobs <- 2
	close(jobs)
	for {
		job, ok := <-jobs
		if !ok {
			break
		}
		fmt.Println(job)
	}
}
\end{lstlisting}
\subsection*{Output}
\begin{lstlisting}[style=terminal]
1
2
\end{lstlisting}
\clearpage
\section{Range over Channels}
\sectionline
range ch receives until the channel is closed. It is the idiomatic way to consume a stream with an explicit end.\par
\subsection*{Example}
\begin{lstlisting}[style=gocode]
package main

import "fmt"

func main() {
	ch := make(chan string, 2)
	ch <- "one"
	ch <- "two"
	close(ch)
	for msg := range ch {
		fmt.Println(msg)
	}
}
\end{lstlisting}
\subsection*{Output}
\begin{lstlisting}[style=terminal]
one
two
\end{lstlisting}
\clearpage
\section{Timers}
\sectionline
time.NewTimer fires once. Stop or drain unused timers when appropriate to avoid unnecessary work.\par
\subsection*{Example}
\begin{lstlisting}[style=gocode]
package main

import (
	"fmt"
	"time"
)

func main() {
	timer := time.NewTimer(10 * time.Millisecond)
	defer timer.Stop()
	<-timer.C
	fmt.Println("timer fired")
}
\end{lstlisting}
\subsection*{Output}
\begin{lstlisting}[style=terminal]
timer fired
\end{lstlisting}
\clearpage
\section{Tickers}
\sectionline
A ticker produces periodic events. Always stop a ticker when you are finished with it.\par
\subsection*{Example}
\begin{lstlisting}[style=gocode]
package main

import (
	"fmt"
	"time"
)

func main() {
	ticker := time.NewTicker(10 * time.Millisecond)
	defer ticker.Stop()
	for i := 0; i < 2; i++ {
		<-ticker.C
		fmt.Println("tick")
	}
}
\end{lstlisting}
\subsection*{Output}
\begin{lstlisting}[style=terminal]
tick
tick
\end{lstlisting}
\clearpage
\section{Worker Pools}
\sectionline
A worker pool caps concurrency by running a fixed number of workers that consume tasks from a shared channel.\par
\subsection*{Example}
\begin{lstlisting}[style=gocode]
package main

import (
	"fmt"
	"sync"
)

func worker(jobs <-chan int, results chan<- int, wg *sync.WaitGroup) {
	defer wg.Done()
	for n := range jobs {
		results <- n * n
	}
}

func main() {
	jobs := make(chan int, 3)
	results := make(chan int, 3)
	var wg sync.WaitGroup
	wg.Add(2)
	for i := 0; i < 2; i++ {
		go worker(jobs, results, &wg)
	}
	for _, n := range []int{2, 3, 4} {
		jobs <- n
	}
	close(jobs)
	wg.Wait()
	close(results)
	for n := range results {
		fmt.Println(n)
	}
}
\end{lstlisting}


\Needspace{6\baselineskip}
\subsection*{Output}
\begin{lstlisting}[style=terminal]
4
9
16
\end{lstlisting}
\emph{Example run.} All three squares are printed, but worker scheduling can change their order.
\clearpage
\section{WaitGroups}
\sectionline
sync.WaitGroup coordinates completion of a known number of goroutines. Call Add before starting work, and Done exactly once per task.\par
\subsection*{Example}
\begin{lstlisting}[style=gocode]
package main

import (
	"fmt"
	"sync"
)

func main() {
	var wg sync.WaitGroup
	for i := 0; i < 3; i++ {
		wg.Add(1)
		go func(n int) {
			defer wg.Done()
			fmt.Print(n, " ")
		}(i)
	}
	wg.Wait()
	fmt.Println()
}
\end{lstlisting}


\Needspace{6\baselineskip}
\subsection*{Output}
\begin{lstlisting}[style=terminal]
2 0 1
\end{lstlisting}
\emph{Example run.} The goroutines may print the numbers in any order, depending on scheduling.
\clearpage
\section{Rate Limiting}
\sectionline
A time.Ticker can pace requests. For bursting workloads, a token-bucket limiter uses a buffered token channel.\par
\subsection*{Example}
\begin{lstlisting}[style=gocode]
package main

import (
	"fmt"
	"time"
)

func main() {
	limiter := time.NewTicker(10 * time.Millisecond)
	defer limiter.Stop()
	for i := 1; i <= 3; i++ {
		<-limiter.C
		fmt.Println("request", i)
	}
}
\end{lstlisting}
\subsection*{Output}
\begin{lstlisting}[style=terminal]
request 1
request 2
request 3
\end{lstlisting}
\clearpage
\section{Atomic Counters}
\sectionline
Atomic variables provide lock-free operations on individual values. Use sync/atomic types instead of a shared unsynchronized integer.\par
\subsection*{Example}
\begin{lstlisting}[style=gocode]
package main

import (
	"fmt"
	"sync"
	"sync/atomic"
)

func main() {
	var count atomic.Int64
	var wg sync.WaitGroup
	for i := 0; i < 100; i++ {
		wg.Add(1)
		go func() {
			defer wg.Done()
			count.Add(1)
		}()
	}
	wg.Wait()
	fmt.Println(count.Load())
}
\end{lstlisting}
\subsection*{Output}
\begin{lstlisting}[style=terminal]
100
\end{lstlisting}
\clearpage
\section{Mutexes}
\sectionline
A mutex protects an invariant spanning one or more variables. Lock around both reading and writing shared state.\par
\subsection*{Example}
\begin{lstlisting}[style=gocode]
package main

import (
	"fmt"
	"sync"
)

type Counter struct {
	mu sync.Mutex
	n  int
}

func (c *Counter) Inc() {
	c.mu.Lock()
	defer c.mu.Unlock()
	c.n++
}

func main() {
	c := &Counter{}
	var wg sync.WaitGroup
	for i := 0; i < 10; i++ {
		wg.Add(1)
		go func() { defer wg.Done(); c.Inc() }()
	}
	wg.Wait()
	fmt.Println(c.n)
}
\end{lstlisting}
\subsection*{Output}
\begin{lstlisting}[style=terminal]
10
\end{lstlisting}
\clearpage
\section{Stateful Goroutines}
\sectionline
Another concurrency model assigns ownership of mutable state to one goroutine; requests arrive via channels.\par
\subsection*{Example}
\begin{lstlisting}[style=gocode]
package main

import "fmt"

type request struct {
	delta int
	reply chan int
}

func main() {
	requests := make(chan request)
	go func() {
		total := 0
		for req := range requests {
			total += req.delta
			req.reply <- total
		}
	}()
	reply := make(chan int)
	requests <- request{delta: 5, reply: reply}
	fmt.Println(<-reply)
	close(requests)
}
\end{lstlisting}
\subsection*{Output}
\begin{lstlisting}[style=terminal]
5
\end{lstlisting}
\clearpage
\section{Sorting}
\sectionline
The slices package provides type-safe sorting and membership functions. Sorting changes the slice in place.\par
\subsection*{Example}
\begin{lstlisting}[style=gocode]
package main

import (
	"fmt"
	"slices"
)

func main() {
	nums := []int{3, 1, 2}
	slices.Sort(nums)
	fmt.Println(nums)
	fmt.Println(slices.IsSorted(nums))
}
\end{lstlisting}
\subsection*{Output}
\begin{lstlisting}[style=terminal]
[1 2 3]
true
\end{lstlisting}
\clearpage
\section{Sorting by Functions}
\sectionline
Use slices.SortFunc and a comparison function to order records by a specific field. cmp.Compare handles most ordered values.\par
\subsection*{Example}
\begin{lstlisting}[style=gocode]
package main

import (
	"cmp"
	"fmt"
	"slices"
)

type Person struct {
	Name string
	Age  int
}

func main() {
	people := []Person{{"Bob", 40}, {"Ada", 36}}
	slices.SortFunc(people, func(a, b Person) int {
		return cmp.Compare(a.Age, b.Age)
	})
	fmt.Println(people)
}
\end{lstlisting}
\subsection*{Output}
\begin{lstlisting}[style=terminal]
[{Ada 36} {Bob 40}]
\end{lstlisting}
\clearpage
\section{Panic}
\sectionline
panic is for unrecoverable programming failures, not routine errors. Deferred functions still run as the stack unwinds.\par
\subsection*{Example}
\begin{lstlisting}[style=gocode]
package main

import "fmt"

func main() {
	defer fmt.Println("cleanup")
	fmt.Println("starting")
	panic("unexpected failure")
}
\end{lstlisting}


\Needspace{6\baselineskip}
\subsection*{Output}
\begin{lstlisting}[style=terminal]
starting
cleanup
panic: unexpected failure
... stack trace ...
\end{lstlisting}
\emph{Abbreviated combined stdout/stderr.} The panic prints a runtime stack trace and exits with a nonzero status; the stack details and source line numbers vary.
\clearpage
\section{Defer}
\sectionline
defer schedules a call for when the surrounding function returns. Deferred calls run in last-in, first-out order.\par
\subsection*{Example}
\begin{lstlisting}[style=gocode]
package main

import "fmt"

func main() {
	defer fmt.Println("first deferred")
	defer fmt.Println("second deferred")
	fmt.Println("body")
}
\end{lstlisting}
\subsection*{Output}
\begin{lstlisting}[style=terminal]
body
second deferred
first deferred
\end{lstlisting}
\clearpage
\section{Recover}
\sectionline
recover can intercept a panic only when called from a deferred function in the same goroutine. It should be used at carefully chosen boundaries.\par
\subsection*{Example}
\begin{lstlisting}[style=gocode]
package main

import "fmt"

func safeCall() {
	defer func() {
		if v := recover(); v != nil {
			fmt.Println("recovered:", v)
		}
	}()
	panic("boom")
}

func main() {
	safeCall()
	fmt.Println("still running")
}
\end{lstlisting}
\subsection*{Output}
\begin{lstlisting}[style=terminal]
recovered: boom
still running
\end{lstlisting}
\clearpage
\section{String Functions}
\sectionline
strings provides helpers for matching, trimming, splitting, and replacing text. Most operations return a new string.\par
\subsection*{Example}
\begin{lstlisting}[style=gocode]
package main

import (
	"fmt"
	"strings"
)

func main() {
	fmt.Println(strings.Contains("golang", "lang"))
	fmt.Println(strings.ToUpper("go"))
	fmt.Println(strings.Join([]string{"a", "b"}, ","))
}
\end{lstlisting}
\subsection*{Output}
\begin{lstlisting}[style=terminal]
true
GO
a,b
\end{lstlisting}
\clearpage
\section{String Formatting}
\sectionline
fmt.Printf formats values with verbs such as \%d (integer), \%s (string), \%v (general), and \%.2f (two decimal places).\par
\subsection*{Example}
\begin{lstlisting}[style=gocode]
package main

import "fmt"

func main() {
	name, count, price := "Ada", 3, 2.5
	fmt.Printf("%s: %d items at $%.2f\n", name, count, price)
	fmt.Printf("%T\n", count)
}
\end{lstlisting}
\subsection*{Output}
\begin{lstlisting}[style=terminal]
Ada: 3 items at $2.50
int
\end{lstlisting}
\clearpage
\section{Text Templates}
\sectionline
text/template separates reusable text layout from data and handles escaping according to its template language (unlike html/template).\par
\subsection*{Example}
\begin{lstlisting}[style=gocode]
package main

import (
	"os"
	"text/template"
)

func main() {
	t := template.Must(template.New("greet").Parse("Hello, {{.Name}}!\n"))
	data := struct{ Name string }{Name: "Ada"}
	if err := t.Execute(os.Stdout, data); err != nil {
		panic(err)
	}
}
\end{lstlisting}
\subsection*{Output}
\begin{lstlisting}[style=terminal]
Hello, Ada!
\end{lstlisting}
\clearpage
\section{Regular Expressions}
\sectionline
The regexp package uses RE2 syntax and provides matching, capture, and replacement operations without backtracking explosions.\par
\subsection*{Example}
\begin{lstlisting}[style=gocode]
package main

import (
	"fmt"
	"regexp"
)

func main() {
	re := regexp.MustCompile(`[a-z]+`)
	fmt.Println(re.FindString("123abc456"))
	fmt.Println(re.ReplaceAllString("one TWO three", "x"))
}
\end{lstlisting}
\subsection*{Output}
\begin{lstlisting}[style=terminal]
abc
x TWO x
\end{lstlisting}
\clearpage
\section{JSON}
\sectionline
encoding/json marshals exported struct fields. Decode into a typed struct when the schema is known and handle errors explicitly.\par
\subsection*{Example}
\begin{lstlisting}[style=gocode]
package main

import (
	"encoding/json"
	"fmt"
)

type Person struct {
	Name string `json:"name"`
	Age  int    `json:"age"`
}

func main() {
	b, err := json.Marshal(Person{Name: "Ada", Age: 36})
	if err != nil {
		panic(err)
	}
	fmt.Println(string(b))
	var p Person
	if err := json.Unmarshal(b, &p); err != nil {
		panic(err)
	}
	fmt.Println(p.Name)
}
\end{lstlisting}
\subsection*{Output}
\begin{lstlisting}[style=terminal]
{"name":"Ada","age":36}
Ada
\end{lstlisting}
\clearpage
\section{XML}
\sectionline
encoding/xml maps XML elements and attributes to Go structs. Use field tags when XML element names differ from Go field names.\par
\subsection*{Example}
\begin{lstlisting}[style=gocode]
package main

import (
	"encoding/xml"
	"fmt"
)

type Item struct {
	XMLName xml.Name `xml:"item"`
	ID      int      `xml:"id,attr"`
	Title   string   `xml:"title"`
}

func main() {
	b, err := xml.Marshal(Item{ID: 7, Title: "Go"})
	if err != nil {
		panic(err)
	}
	fmt.Println(string(b))
}
\end{lstlisting}
\subsection*{Output}
\begin{lstlisting}[style=terminal]
<item id="7"><title>Go</title></item>
\end{lstlisting}
\clearpage
\section{Time}
\sectionline
time.Time represents an instant; time.Duration represents elapsed time. Prefer Add and Sub for arithmetic.\par
\subsection*{Example}
\begin{lstlisting}[style=gocode]
package main

import (
	"fmt"
	"time"
)

func main() {
	start := time.Date(2026, 1, 2, 3, 4, 5, 0, time.UTC)
	end := start.Add(90 * time.Minute)
	fmt.Println(end.Format(time.RFC3339))
	fmt.Println(end.Sub(start))
}
\end{lstlisting}
\subsection*{Output}
\begin{lstlisting}[style=terminal]
2026-01-02T04:34:05Z
1h30m0s
\end{lstlisting}
\clearpage
\section{Epoch}
\sectionline
Unix timestamps count elapsed seconds or smaller units from 1970-01-01 UTC. Do not confuse seconds, milliseconds, and nanoseconds.\par
\subsection*{Example}
\begin{lstlisting}[style=gocode]
package main

import (
	"fmt"
	"time"
)

func main() {
	t := time.Unix(0, 0).UTC()
	fmt.Println(t.Format(time.RFC3339))
	fmt.Println(t.Unix())
}
\end{lstlisting}
\subsection*{Output}
\begin{lstlisting}[style=terminal]
1970-01-01T00:00:00Z
0
\end{lstlisting}
\clearpage
\section{Time Formatting / Parsing}
\sectionline
Go uses a reference timestamp rather than printf-style placeholders: 2006-01-02 15:04:05.\par
\subsection*{Example}
\begin{lstlisting}[style=gocode]
package main

import (
	"fmt"
	"time"
)

func main() {
	t, err := time.Parse("2006-01-02", "2026-10-21")
	if err != nil {
		panic(err)
	}
	fmt.Println(t.Format("Jan 2, 2006"))
	fmt.Println(t.Format(time.RFC3339))
}
\end{lstlisting}
\subsection*{Output}
\begin{lstlisting}[style=terminal]
Oct 21, 2026
2026-10-21T00:00:00Z
\end{lstlisting}
\clearpage
\section{Random Numbers}
\sectionline
math/rand/v2 generates pseudorandom values for simulations and sampling. Do not use it for secrets; use crypto/rand instead.\par
\subsection*{Example}
\begin{lstlisting}[style=gocode]
package main

import (
	"fmt"
	"math/rand/v2"
)

func main() {
	n := rand.IntN(10) // 0 through 9
	fmt.Println(n >= 0 && n < 10)
}
\end{lstlisting}
\subsection*{Output}
\begin{lstlisting}[style=terminal]
true
\end{lstlisting}
\clearpage
\section{Number Parsing}
\sectionline
strconv converts strings into numeric types and returns errors when input is invalid.\par
\subsection*{Example}
\begin{lstlisting}[style=gocode]
package main

import (
	"fmt"
	"strconv"
)

func main() {
	n, err := strconv.Atoi("42")
	if err != nil {
		panic(err)
	}
	fmt.Println(n + 1)
	fmt.Println(strconv.FormatBool(true))
}
\end{lstlisting}
\subsection*{Output}
\begin{lstlisting}[style=terminal]
43
true
\end{lstlisting}
\clearpage
\section{URL Parsing}
\sectionline
net/url parses URL components, and url.Values encodes query parameters. Never concatenate untrusted query values by hand.\par
\subsection*{Example}
\begin{lstlisting}[style=gocode]
package main

import (
	"fmt"
	"net/url"
)

func main() {
	u, err := url.Parse("https://example.com/search?q=go%20lang")
	if err != nil {
		panic(err)
	}
	fmt.Println(u.Host)
	fmt.Println(u.Query().Get("q"))
}
\end{lstlisting}
\subsection*{Output}
\begin{lstlisting}[style=terminal]
example.com
go lang
\end{lstlisting}
\clearpage
\section{SHA256 Hashes}
\sectionline
SHA-256 maps bytes to a fixed-size digest. It is a hash, not encryption; do not use plain SHA-256 to store passwords.\par
\subsection*{Example}
\begin{lstlisting}[style=gocode]
package main

import (
	"crypto/sha256"
	"fmt"
)

func main() {
	digest := sha256.Sum256([]byte("hello"))
	fmt.Printf("%x\n", digest)
}
\end{lstlisting}
\subsection*{Output}
\begin{lstlisting}[style=terminal]
2cf24dba5fb0a30e26e83b2ac5b9e29e1b161e5c1fa7425e73043362938b9824
\end{lstlisting}
\clearpage
\section{Base64 Encoding}
\sectionline
Base64 transforms bytes into printable text; it does not provide confidentiality. Choose the URL-safe encoding for URL components.\par
\subsection*{Example}
\begin{lstlisting}[style=gocode]
package main

import (
	"encoding/base64"
	"fmt"
)

func main() {
	encoded := base64.StdEncoding.EncodeToString([]byte("hello"))
	decoded, err := base64.StdEncoding.DecodeString(encoded)
	if err != nil {
		panic(err)
	}
	fmt.Println(encoded)
	fmt.Println(string(decoded))
}
\end{lstlisting}
\subsection*{Output}
\begin{lstlisting}[style=terminal]
aGVsbG8=
hello
\end{lstlisting}
\clearpage
\section{Reading Files}
\sectionline
os.ReadFile is convenient for small files. For streaming or large files, prefer bufio.Scanner or buffered readers.\par
\subsection*{Example}
\begin{lstlisting}[style=gocode]
package main

import (
	"fmt"
	"os"
)

func main() {
	path := "example.txt"
	if err := os.WriteFile(path, []byte("hello\n"), 0600); err != nil {
		panic(err)
	}
	defer os.Remove(path)
	data, err := os.ReadFile(path)
	if err != nil {
		panic(err)
	}
	fmt.Print(string(data))
}
\end{lstlisting}
\subsection*{Output}
\begin{lstlisting}[style=terminal]
hello
\end{lstlisting}
\clearpage
\section{Writing Files}
\sectionline
os.WriteFile atomically replaces neither the file nor its directory entry: for crash-safe updates, write a temporary file and rename it.\par
\subsection*{Example}
\begin{lstlisting}[style=gocode]
package main

import (
	"fmt"
	"os"
)

func main() {
	path := "out.txt"
	err := os.WriteFile(path, []byte("saved\n"), 0600)
	if err != nil {
		panic(err)
	}
	defer os.Remove(path)
	fmt.Println("wrote", path)
}
\end{lstlisting}
\subsection*{Output}
\begin{lstlisting}[style=terminal]
wrote out.txt
\end{lstlisting}
\clearpage
\section{Line Filters}
\sectionline
A bufio.Scanner reads input line by line. Check scanner.Err after iteration; increase the buffer for very long lines.\par
\subsection*{Example}
\begin{lstlisting}[style=gocode]
package main

import (
	"bufio"
	"fmt"
	"strings"
)

func main() {
	scanner := bufio.NewScanner(strings.NewReader("one\ntwo\n"))
	for scanner.Scan() {
		fmt.Println(strings.ToUpper(scanner.Text()))
	}
	if err := scanner.Err(); err != nil {
		panic(err)
	}
}
\end{lstlisting}
\subsection*{Output}
\begin{lstlisting}[style=terminal]
ONE
TWO
\end{lstlisting}
\clearpage
\section{File Paths}
\sectionline
filepath.Join builds OS-native paths; filepath.Base and filepath.Ext inspect components without manually splitting separators.\par
\subsection*{Example}
\begin{lstlisting}[style=gocode]
package main

import (
	"fmt"
	"path/filepath"
)

func main() {
	path := filepath.Join("docs", "notes.txt")
	fmt.Println(path)
	fmt.Println(filepath.Base(path))
	fmt.Println(filepath.Ext(path))
}
\end{lstlisting}
\subsection*{Output}
\begin{lstlisting}[style=terminal]
docs/notes.txt
notes.txt
.txt
\end{lstlisting}
\clearpage
\section{Directories}
\sectionline
os.MkdirAll creates nested directories. os.ReadDir lists entries; filepath.WalkDir visits a whole tree with error handling.\par
\subsection*{Example}
\begin{lstlisting}[style=gocode]
package main

import (
	"fmt"
	"os"
	"path/filepath"
)

func main() {
	root, err := os.MkdirTemp("", "tree-*")
	if err != nil {
		panic(err)
	}
	defer os.RemoveAll(root)
	path := filepath.Join(root, "a", "b")
	if err := os.MkdirAll(path, 0755); err != nil {
		panic(err)
	}
	entries, err := os.ReadDir(filepath.Join(root, "a"))
	if err != nil {
		panic(err)
	}
	for _, e := range entries {
		fmt.Println(e.Name())
	}
}
\end{lstlisting}
\subsection*{Output}
\begin{lstlisting}[style=terminal]
b
\end{lstlisting}
\clearpage
\section{Temporary Files and Directories}
\sectionline
os.CreateTemp and os.MkdirTemp choose unique names safely. Remove temporary files and directories when no longer needed.\par
\subsection*{Example}
\begin{lstlisting}[style=gocode]
package main

import (
	"fmt"
	"os"
)

func main() {
	f, err := os.CreateTemp("", "demo-*.txt")
	if err != nil {
		panic(err)
	}
	defer os.Remove(f.Name())
	defer f.Close()
	fmt.Println("temporary file created:", f.Name() != "")
}
\end{lstlisting}
\subsection*{Output}
\begin{lstlisting}[style=terminal]
temporary file created: true
\end{lstlisting}
\clearpage
\section{Embed Directive}
\sectionline
The //go:embed directive includes a file at compile time. The named file must exist relative to the package directory.\par
\subsection*{Example}
\begin{lstlisting}[style=gocode]
package main

import (
	_ "embed"
	"fmt"
)

//go:embed message.txt
var message string

func main() {
	fmt.Print(message)
}
\end{lstlisting}
Create a file named \texttt{message.txt} in the same directory containing \texttt{hello from a file} followed by a newline before building the program.
\subsection*{Output}
\begin{lstlisting}[style=terminal]
hello from a file
\end{lstlisting}
\clearpage
\section{Testing and Benchmarking}
\sectionline
Tests live in files ending with \_test.go. A test function receives *testing.T; benchmarks receive *testing.B.\par
\subsection*{Example: calc.go}
\begin{lstlisting}[style=gocode]
package calc

func Add(a, b int) int { return a + b }
\end{lstlisting}
Put the function above in \texttt{calc.go} and the following test in \texttt{calc\_test.go}:
\subsection*{Example: calc\_test.go}
\begin{lstlisting}[style=gocode]
package calc

import "testing"

func TestAdd(t *testing.T) {
    if got := Add(2, 3); got != 5 {
        t.Fatalf("Add(2, 3) = %d; want 5", got)
    }
}

func BenchmarkAdd(b *testing.B) {
    for i := 0; i < b.N; i++ {
        Add(2, 3)
    }
}
\end{lstlisting}


\Needspace{6\baselineskip}
\subsection*{Output}
\begin{lstlisting}[style=terminal]
=== RUN   TestAdd
--- PASS: TestAdd (0.00s)
BenchmarkAdd-5  189387514  1.841 ns/op
PASS
ok  example/calc  0.453s
\end{lstlisting}
\emph{Example using \texttt{go test -v -bench=.} (tests and benchmarks).} Benchmark iterations, nanoseconds per operation, CPU-dependent suffixes, package path, and elapsed times vary. Tests report \texttt{PASS} when assertions succeed.
\clearpage
\section{Command-Line Arguments}
\sectionline
os.Args contains the executable path as element zero and then command-line arguments. Validate the count before indexing.\par
\subsection*{Example}
\begin{lstlisting}[style=gocode]
package main

import (
	"fmt"
	"os"
)

func main() {
	if len(os.Args) < 2 {
		fmt.Println("usage: app NAME")
		return
	}
	fmt.Println("hello", os.Args[1])
}
\end{lstlisting}


\Needspace{6\baselineskip}
\subsection*{Output}
\begin{lstlisting}[style=terminal]
hello Ada
\end{lstlisting}
\emph{Example using \texttt{go run main.go Ada}.} The greeting changes with the first argument. With no argument, the program prints usage information instead.
\clearpage
\section{Command-Line Flags}
\sectionline
The flag package parses typed command-line options. Define flags before calling flag.Parse.\par
\subsection*{Example}
\begin{lstlisting}[style=gocode]
package main

import (
	"flag"
	"fmt"
)

func main() {
	name := flag.String("name", "world", "name to greet")
	flag.Parse()
	fmt.Println("hello", *name)
}
\end{lstlisting}


\Needspace{6\baselineskip}
\subsection*{Output}
\begin{lstlisting}[style=terminal]
hello Ada
\end{lstlisting}
\emph{Example using \texttt{go run main.go -name Ada}.} The flag controls the name; omitting it prints \texttt{hello world}.
\clearpage
\section{Command-Line Subcommands}
\sectionline
Subcommands can use separate flag.FlagSet values. Check the command first, then parse only its arguments.\par
\subsection*{Example}
\begin{lstlisting}[style=gocode]
package main

import (
	"flag"
	"fmt"
	"os"
)

func main() {
	if len(os.Args) < 2 || os.Args[1] != "greet" {
		fmt.Println("usage: app greet [-name NAME]")
		return
	}
	greet := flag.NewFlagSet("greet", flag.ExitOnError)
	name := greet.String("name", "world", "name")
	greet.Parse(os.Args[2:])
	fmt.Println("hello", *name)
}
\end{lstlisting}


\Needspace{6\baselineskip}
\subsection*{Output}
\begin{lstlisting}[style=terminal]
hello Ada
\end{lstlisting}
\emph{Example using \texttt{go run main.go greet -name Ada}.} The output depends on the subcommand and flags; invalid or missing subcommands print usage information.
\clearpage
\section{Environment Variables}
\sectionline
Use os.LookupEnv when you need to distinguish an unset environment variable from an empty one.\par
\subsection*{Example}
\begin{lstlisting}[style=gocode]
package main

import (
	"fmt"
	"os"
)

func main() {
	os.Setenv("APP_MODE", "test")
	mode, ok := os.LookupEnv("APP_MODE")
	fmt.Println(mode, ok)
}
\end{lstlisting}
\subsection*{Output}
\begin{lstlisting}[style=terminal]
test true
\end{lstlisting}
\clearpage
\section{Logging}
\sectionline
log/slog provides structured logs with named fields and pluggable handlers. Avoid secrets in log output.\par
\subsection*{Example}
\begin{lstlisting}[style=gocode]
package main

import (
	"log/slog"
	"os"
)

func main() {
	log := slog.New(slog.NewTextHandler(os.Stdout, nil))
	log.Info("request handled", "method", "GET", "status", 200)
}
\end{lstlisting}


\Needspace{6\baselineskip}
\subsection*{Output}
\begin{lstlisting}[style=terminal]
time=2026-10-21T14:21:00Z level=INFO msg="request handled" method=GET status=200
\end{lstlisting}
\emph{Example run.} The timestamp changes for each invocation. Log fields and formatting depend on the configured handler.
\clearpage
\section{HTTP Client}
\sectionline
http.Client makes requests. Always close response bodies and check status codes before treating a response as successful.\par
\subsection*{Example}
\begin{lstlisting}[style=gocode]
package main

import (
	"fmt"
	"net/http"
	"time"
)

func main() {
	client := &http.Client{Timeout: 5 * time.Second}
	resp, err := client.Get("https://example.com")
	if err != nil {
		panic(err)
	}
	defer resp.Body.Close()
	fmt.Println("status:", resp.StatusCode)
}
\end{lstlisting}


\Needspace{6\baselineskip}
\subsection*{Output}
\begin{lstlisting}[style=terminal]
status: 200
\end{lstlisting}
\emph{Example response.} The status code depends on the remote server; DNS, TLS, timeouts, or connectivity can cause the request to fail instead.
\clearpage
\section{HTTP Server}
\sectionline
An HTTP handler writes a response through http.ResponseWriter. ServeMux routes method/path patterns to handlers (Go 1.22+).\par
\subsection*{Example}
\begin{lstlisting}[style=gocode]
package main

import (
	"fmt"
	"log"
	"net/http"
)

func main() {
	mux := http.NewServeMux()
	mux.HandleFunc("GET /hello", func(w http.ResponseWriter, r *http.Request) {
		fmt.Fprintln(w, "hello")
	})
	log.Fatal(http.ListenAndServe(":8080", mux))
}
\end{lstlisting}


\Needspace{6\baselineskip}
\subsection*{Output}
\begin{lstlisting}[style=terminal]
hello
\end{lstlisting}
\emph{Example HTTP response body.} Run the server, then request \texttt{GET /hello} on port 8080 from another terminal or browser. The server keeps running; it does not print this body to its own terminal.
\clearpage
\section{TCP Server}
\sectionline
net.Listen accepts incoming TCP connections. Handle each client separately and close connections when finished.\par
\subsection*{Example}
\begin{lstlisting}[style=gocode]
package main

import (
	"bufio"
	"fmt"
	"log"
	"net"
)

func main() {
	ln, err := net.Listen("tcp", ":8080")
	if err != nil {
		log.Fatal(err)
	}
	defer ln.Close()
	for {
		conn, err := ln.Accept()
		if err != nil {
			log.Print(err)
			continue
		}
		go func(c net.Conn) {
			defer c.Close()
			text, err := bufio.NewReader(c).ReadString('\n')
			if err == nil {
				fmt.Fprint(c, text)
			}
		}(conn)
	}
}
\end{lstlisting}


\Needspace{6\baselineskip}
\subsection*{Output}
\begin{lstlisting}[style=terminal]
ping
\end{lstlisting}
\emph{Example client response.} After connecting to port 8080, a client sending \texttt{ping} followed by a newline receives that line back. The server runs until stopped; only the first line of each connection is echoed.
\clearpage
\section{Context}
\sectionline
context.Context carries cancellation and deadlines through a call chain. Pass it as the first argument rather than storing it on a struct.\par
\subsection*{Example}
\begin{lstlisting}[style=gocode]
package main

import (
	"context"
	"fmt"
	"time"
)

func main() {
	ctx, cancel := context.WithCancel(context.Background())
	cancel()
	select {
	case <-ctx.Done():
		fmt.Println(ctx.Err())
	case <-time.After(time.Second):
		fmt.Println("finished")
	}
}
\end{lstlisting}
\subsection*{Output}
\begin{lstlisting}[style=terminal]
context canceled
\end{lstlisting}
\clearpage
\section{Spawning Processes}
\sectionline
os/exec starts processes without invoking a shell by default. Avoid constructing shell command strings from user input.\par
\subsection*{Example}
\begin{lstlisting}[style=gocode]
package main

import (
	"fmt"
	"os/exec"
)

func main() {
	out, err := exec.Command("go", "version").Output()
	if err != nil {
		panic(err)
	}
	fmt.Print(string(out))
}
\end{lstlisting}


\Needspace{6\baselineskip}
\subsection*{Output}
\begin{lstlisting}[style=terminal]
go version go1.23.2 linux/amd64
\end{lstlisting}
\emph{Example run.} The Go version, operating system, and architecture match the toolchain installed on the host and therefore vary.
\clearpage
\section{Exec'ing Processes}
\sectionline
syscall.Exec replaces the current process on Unix. It does not return on success and is generally inappropriate for portable applications.\par
\subsection*{Example}
\begin{lstlisting}[style=gocode]
package main

import (
	"os"
	"os/exec"
	"syscall"
)

func main() {
	binary, err := exec.LookPath("ls")
	if err != nil {
		panic(err)
	}
	if err := syscall.Exec(binary, []string{"ls", "-l"}, os.Environ()); err != nil {
		panic(err)
	}
}
\end{lstlisting}


\Needspace{6\baselineskip}
\subsection*{Output}
\begin{lstlisting}[style=terminal]
total 4
-rw-r--r-- 1 user user 237 Oct 21  main.go
\end{lstlisting}
\emph{Example output on Unix.} The process becomes \texttt{ls -l}; directory contents, owners, sizes, permissions, and dates depend on the working directory and operating system.
\clearpage
\section{Signals}
\sectionline
signal.NotifyContext cancels a context when one of the specified operating-system signals arrives. stop releases signal resources.\par
\subsection*{Example}
\begin{lstlisting}[style=gocode]
package main

import (
	"context"
	"fmt"
	"os/signal"
	"syscall"
)

func main() {
	ctx, stop := signal.NotifyContext(
		context.Background(),
		syscall.SIGINT, syscall.SIGTERM,
	)
	defer stop()

	fmt.Println("waiting for signal")
	<-ctx.Done()
	fmt.Println("shutting down")
}
\end{lstlisting}


\Needspace{6\baselineskip}
\subsection*{Output}
\begin{lstlisting}[style=terminal]
waiting for signal
shutting down
\end{lstlisting}
\emph{Example after Ctrl+C or SIGTERM.} The second line appears only after a signal arrives. How long the program waits depends on when that happens.
\clearpage
\section{Exit}
\sectionline
os.Exit terminates immediately with a status code and does not execute deferred calls. Returning from main is preferable for routine success.\par
\subsection*{Example}
\begin{lstlisting}[style=gocode]
package main

import (
	"fmt"
	"os"
)

func main() {
	fmt.Fprintln(os.Stderr, "invalid arguments")
	os.Exit(2)
}
\end{lstlisting}


\Needspace{6\baselineskip}
\subsection*{Output}
\begin{lstlisting}[style=terminal]
invalid arguments
\end{lstlisting}
\emph{Example stderr output.} The process exits with status code \texttt{2}, which a calling shell can inspect. Deferred calls do not run when \texttt{os.Exit} executes.
\end{document}
